% Section 11.3, from lectures/L11/README.md: what you should be able to do, the questions to test
% yourself with, and the next lecture, with the README's advice to prepare for the exam with the
% practice exam.

\needspace{12\baselineskip}
\section{Sammanfattning}
\label{c11:sec:summary}

\needspace{6\baselineskip}
\subsection*{Det här ska du kunna nu}
Efter det här kapitlet ska du kunna:
\begin{itemize}
  \item Beskriva ett styrsystem som ingångar, beslut och utgångar, och peka ut dem i ett exempel.
  \item Räkna ut förkopplingsmotståndet för en lysdiod, och förklara varför det behövs.
  \item Förklara hur en transistor och ett relä låter mikrodatorn styra en last på 24~V, och vad
    dioden över spolen gör.
  \item Förklara varför en knapp studsar, när det spelar roll, och hur en kort väntan löser det.
  \item Gå från en tillståndstabell till en flödesplan och ett program, och testa programmet först
    i simulatorn och sedan på ett riktigt kort.
  \item Felsöka ett program på kortet: skilja ett programfel från ett kopplingsfel.
\end{itemize}

\needspace{6\baselineskip}
\subsection*{Frågor att testa dig själv med}
\begin{enumerate}
  \item Varför får en lysdiod aldrig kopplas direkt mellan ett stift och jord?
  \item Ett stift tål 40~mA. Varför styrs ett relä ändå genom en transistor?
  \item Ett program som räknar knapptryck räknar för många, men bara på kortet, aldrig i
    simulatorn. Varför?
  \item Vad händer om \code{QUARTER_MS} står kvar på 1 när programmet förs över till kortet?
  \item Programmet räknar en sekund exakt på cykeln. Varför kan trafikljuset ändå gå fel med en
    halv procent?
  \item Varför skriver \code{out PORTB, r16} hela porten på en gång, och varför är det en fördel i
    ett trafikljus?
\end{enumerate}

\needspace{6\baselineskip}
\subsection*{Nästa kapitel}
Kapitel~\ref{c12:ch} förbereder det skriftliga provet:
\begin{itemize}
  \item Repetition av hela kursen, lärandemål för lärandemål.
  \item Det skriftliga provet: två timmar, utan miniräknare, med referensbladet för
    AVR-instruktionerna (bilaga~\ref{app:avr}).
\end{itemize}
Förbered provet med övningsprovet i bilaga~\ref{app:practice}.
